\chapter{Estado del arte}
\label{chap:estado_arte}

\lettrine{E}{l} desarrollo de sistemas \acrfull{bci} ha experimentado un crecimiento exponencial en las últimas décadas, impulsado por avances significativos en neurociencia computacional, procesamiento de señales digitales y algoritmos de \acrfull{ml}. Este crecimiento ha dado lugar a un ecosistema diverso de herramientas, dispositivos y metodologías que abordan diferentes aspectos del desarrollo \acrshort{bci}, desde la adquisición de señales cerebrales hasta el análisis de datos y la implementación de aplicaciones prácticas.

La diversidad del panorama actual refleja tanto la riqueza del campo como su fragmentación inherente. Diferentes grupos de investigación y empresas han desarrollado soluciones especializadas que, aunque robustas en sus dominios específicos, presentan limitaciones cuando se requiere una integración completa para el desarrollo de aplicaciones \acrshort{bci} efectivas. Esta fragmentación constituye uno de los principales obstáculos para la democratización y adopción masiva de tecnologías \acrshort{bci}.

El presente capítulo analiza de forma exhaustiva el panorama actual de herramientas y tecnologías disponibles para el desarrollo de sistemas \acrshort{bci}, identificando sus fortalezas, limitaciones y las oportunidades que estas limitaciones revelan para el desarrollo de soluciones más integradas y accesibles.

\section{Interfaces cerebro-computadora: panorama tecnológico actual}
\label{sec:bci_panorama}

Los sistemas \acrshort{bci} contemporáneos representan la convergencia de múltiples disciplinas científicas y tecnológicas, incluyendo neurociencia, ingeniería biomédica, procesamiento de señales, inteligencia artificial y diseño de interfaces humano-computadora. Esta convergencia ha resultado en una diversidad de enfoques técnicos y metodológicos que reflejan las diferentes prioridades y contextos de aplicación de los sistemas \acrshort{bci}.

La evolución histórica del campo \acrshort{bci} ha estado marcada por hitos tecnológicos significativos que han ampliado progresivamente las posibilidades de aplicación. Desde los primeros experimentos de control de cursores mediante señales \acrshort{eeg} hasta las modernas interfaces híbridas que combinan múltiples modalidades de señales cerebrales, el campo ha demostrado una capacidad notable de innovación y adaptación a nuevos desafíos técnicos y aplicaciones prácticas.

La maduración tecnológica del campo se refleja en la disponibilidad creciente de dispositivos comerciales accesibles, el desarrollo de estándares de interoperabilidad y la emergencia de comunidades de código abierto que facilitan la colaboración y el intercambio de conocimientos. Sin embargo, esta maduración también ha revelado nuevos desafíos relacionados con la integración de herramientas, la estandarización de metodologías y la necesidad de soluciones más accesibles para investigadores con diferentes niveles de experiencia técnica.

\subsection{Técnicas de adquisición de señales cerebrales}
\label{subsec:tecnicas_adquisicion}

La selección de la modalidad de adquisición de señales cerebrales constituye una de las decisiones más fundamentales en el diseño de sistemas \acrshort{bci}, ya que determina tanto las capacidades como las limitaciones del sistema resultante. Cada técnica presenta características distintivas en términos de resolución temporal y espacial, invasividad, coste, complejidad de implementación y tipos de señales cerebrales que puede capturar de forma efectiva.

La \acrfull{eeg} se ha establecido como la modalidad predominante para aplicaciones \acrshort{bci} no invasivas debido a su excepcional resolución temporal, que permite capturar cambios en la actividad eléctrica cerebral con precisión de milisegundos~\cite{eeg_bci_review,wolpaw2002brain}. Esta técnica registra las señales eléctricas generadas por la actividad sincronizada de poblaciones neuronales mediante electrodos colocados en el cuero cabelludo siguiendo sistemas estandarizados como el 10-20 o sus extensiones modernas. La naturaleza no invasiva de la \acrshort{eeg}, combinada con la disponibilidad creciente de dispositivos portátiles y de bajo coste, la convierte en la opción más accesible para investigación y aplicaciones prácticas.

La \acrfull{meg} representa una alternativa sofisticada que detecta los campos magnéticos sutiles producidos por la actividad neuronal, ofreciendo una resolución espacial superior a la \acrshort{eeg} mientras mantiene una excelente resolución temporal~\cite{eeg_bci_review}. Sin embargo, los sistemas \acrshort{meg} requieren entornos de blindaje magnético especializados y equipamiento extremadamente costoso, limitando su aplicabilidad principalmente a centros de investigación avanzados y aplicaciones médicas especializadas.

La \acrfull{fnirs} emplea espectroscopía de infrarrojo cercano para medir cambios en la oxigenación sanguínea cerebral, proporcionando una ventana indirecta pero valiosa hacia la actividad neuronal~\cite{eeg_bci_review}. Esta técnica presenta ventajas significativas en términos de robustez ante artefactos de movimiento y facilidad de implementación, aunque su resolución temporal limitada y la naturaleza indirecta de las mediciones restringen su aplicabilidad a ciertos tipos de paradigmas \acrshort{bci}.

Las técnicas invasivas, incluyendo la \acrfull{ecog} y los registros intracraneales mediante microelectrodos, ofrecen la mejor combinación de resolución temporal y espacial disponible actualmente~\cite{dornhege2007toward}. Estos métodos proporcionan acceso directo a señales neuronales de alta fidelidad y permiten el control preciso de dispositivos externos, pero conllevan riesgos quirúrgicos significativos que limitan su aplicación a casos médicos específicos donde los beneficios potenciales justifican los riesgos asociados.

\section{Dispositivos de adquisición \acrshort{eeg}}
\label{sec:dispositivos_eeg}

El mercado de dispositivos \acrshort{eeg} ha experimentado una diversificación notable en la última década, con productos que abarcan desde sistemas clínicos de precisión extrema hasta dispositivos de consumo orientados a aplicaciones de bienestar y entretenimiento. Esta diversificación refleja tanto la maduración tecnológica del campo como la emergencia de nuevos mercados y aplicaciones que demandan diferentes combinaciones de precisión, portabilidad, facilidad de uso y coste.

La segmentación del mercado se ha establecido naturalmente en tres categorías principales que atienden a diferentes necesidades y contextos de aplicación. Esta segmentación no es rígida, y muchos dispositivos presentan características que los posicionan en zonas intermedias entre categorías, reflejando la evolución continua de las tecnologías subyacentes y las demandas cambiantes del mercado.

\subsection{Sistemas clínicos de alta gama}
\label{subsec:sistemas_clinicos}

Los dispositivos de esta categoría se caracterizan por ofrecer la máxima precisión y capacidades de adquisición disponibles, diseñados específicamente para aplicaciones médicas y de investigación avanzada donde la calidad de señal y la confiabilidad son críticas. Estos sistemas típicamente incorporan entre 64 y 256 canales con resolución de 24 bits y frecuencias de muestreo que pueden alcanzar los 20 kHz, proporcionando una fidelidad de señal excepcional que permite análisis sofisticados de la actividad cerebral.

Los sistemas de \textit{Neuroscan}~\cite{neuroscan_systems} representan el estándar de oro en muchos laboratorios de neurociencia, ofreciendo configuraciones modulares que pueden adaptarse a necesidades específicas de investigación. Los productos de \textit{Brain Products}~\cite{brain_products_systems}, incluyendo la línea \textit{actiCAP}, han ganado reconocimiento por su excelente calidad de señal y facilidad de uso en entornos de investigación. \textit{ANT Neuro}~\cite{ant_neuro_systems} se ha especializado en sistemas de alta densidad que permiten estudios detallados de localización de fuentes y conectividad cerebral.

Estos sistemas incorporan características avanzadas como impedancia de entrada superior a 100 MΩ, rechazo de modo común excepcional y capacidades de sincronización precisas con otros equipos de laboratorio. Sin embargo, su coste elevado, que típicamente oscila entre €50,000 y €200,000, y la complejidad de su operación los limita a instituciones con recursos significativos y personal técnico especializado.

\subsection{Sistemas de investigación de gama media}
\label{subsec:sistemas_investigacion}

Esta categoría representa un equilibrio cuidadosamente optimizado entre capacidades técnicas, facilidad de uso y coste, dirigiéndose específicamente a laboratorios de investigación, instituciones educativas y desarrolladores de aplicaciones \acrshort{bci} que requieren calidad profesional sin el coste prohibitivo de los sistemas clínicos de alta gama.

El \textit{g.tec Unicorn Hybrid Black}~\cite{gtec_unicorn} ha emergido como un referente en esta categoría, ofreciendo 8 canales \acrshort{eeg} de alta calidad con resolución de 24 bits y frecuencia de muestreo de 250 Hz, complementados con sensores inerciales integrados. Su diseño híbrido y conectividad \textit{Bluetooth} de baja latencia lo hacen especialmente atractivo para aplicaciones móviles y estudios en entornos naturales.

El \textit{Emotiv EPOC X}~\cite{emotiv_epoc_x} proporciona 14 canales con opciones de resolución configurable y ha ganado popularidad por su facilidad de configuración y soporte de software extensivo. El \textit{OpenBCI Cyton}~\cite{openbci_cyton} destaca por su naturaleza de código abierto y modularidad, permitiendo configuraciones personalizadas que pueden adaptarse a necesidades específicas de investigación.

Los dispositivos de esta categoría típicamente ofrecen entre 8 y 32 canales con resoluciones de 16 a 24 bits y frecuencias de muestreo de 250 a 1000 Hz. La conectividad inalámbrica es una característica común, facilitando estudios de movilidad y reduciendo las restricciones de cableado. Los costes oscilan entre €1,000 y €10,000, posicionándolos como opciones accesibles para una amplia gama de aplicaciones de investigación.

\begin{table}[hp!]
\centering
\resizebox{\textwidth}{!}{
\begin{tabular}{l|c|c|c|c|c}
\rowcolor{udcpink!25}
\textbf{Dispositivo} & \textbf{Canales} & \textbf{Frecuencia} & \textbf{Resolución} & \textbf{Precio} & \textbf{Consultado} \\
\hline
\textbf{g.tec Unicorn}~\cite{gtec_unicorn} & 8 & 250 Hz & 24 bits & €1,200 & Feb 2025 \\
\textbf{Emotiv EPOC X}~\cite{emotiv_epoc_x} & 14 & 128/256 Hz & 14/16 bits & €999 & Feb 2025 \\
\textbf{OpenBCI Cyton}~\cite{openbci_cyton} & 8 & 250 Hz & 24 bits & €947 & Feb 2025 \\
\textbf{Neurosity Crown}~\cite{neurosity_crown} & 8 & 256 Hz & 16 bits & €899 & Feb 2025 \\
\textbf{Neuroelectrics Enobio}~\cite{neuroelectrics_enobio} & 32 & 500 Hz & 24 bits & €5,000 & Feb 2025 \\
\end{tabular}
}
\caption{Comparativa de dispositivos \acrshort{eeg} de investigación de gama media}
\label{tab:dispositivos_gama_media}
\end{table}

\subsection{Dispositivos de consumo}
\label{subsec:dispositivos_consumo}

Los dispositivos de esta categoría han sido diseñados específicamente para democratizar el acceso a tecnologías \acrshort{eeg}, priorizando la facilidad de uso, portabilidad y accesibilidad económica sobre la precisión técnica extrema. Estos productos han abierto nuevos mercados en aplicaciones de bienestar, meditación, entrenamiento cognitivo y entretenimiento interactivo.

El \textit{InteraXon Muse 2}~\cite{muse_2} se ha posicionado como líder en aplicaciones de meditación y \textit{neurofeedback}, ofreciendo 4 canales \acrshort{eeg} optimizados para la detección de estados de relajación y concentración. Su integración con aplicaciones móviles y ecosistema de contenido lo han convertido en una puerta de entrada popular para usuarios no técnicos interesados en el \textit{neurofeedback}.

El \textit{Emotiv Insight}~\cite{emotiv_insight} proporciona 5 canales con capacidades de detección de estados mentales y expresiones faciales, dirigiéndose a aplicaciones de investigación de mercado y desarrollo de interfaces adaptativas. El \textit{NeuroSky MindWave}~\cite{neurosky_mindwave} representa el extremo más accesible del espectro, ofreciendo funcionalidades básicas de detección de atención y meditación en un factor de forma simple y económico.

Estos dispositivos típicamente incorporan entre 4 y 14 canales con resoluciones de 12 a 16 bits y frecuencias de muestreo de 128 a 512 Hz. El enfoque en facilidad de uso se refleja en configuraciones automatizadas, aplicaciones móviles integradas y documentación orientada a usuarios no técnicos. Los costes oscilan entre €200 y €1,500, posicionándolos como opciones accesibles para exploración personal, educación y desarrollo de prototipos.

\section{Librerías de abstracción de hardware}
\label{sec:librerias_abstraccion}

La diversidad de dispositivos \acrshort{eeg} y sus interfaces propietarias ha motivado el desarrollo de librerías de abstracción que proporcionan \acrfull{api} unificadas para acceder a diferentes tipos de hardware. Estas librerías representan un componente crítico del ecosistema \acrshort{bci}, ya que permiten el desarrollo de aplicaciones que pueden funcionar con múltiples tipos de dispositivos sin requerir modificaciones específicas de hardware.

El desarrollo de estas librerías ha sido impulsado por la necesidad de reducir la fragmentación del ecosistema y facilitar la adopción de tecnologías \acrshort{bci} por parte de desarrolladores que no desean especializarse en los detalles técnicos específicos de cada fabricante de hardware. Sin embargo, la abstracción inevitablemente conlleva compromisos entre simplicidad de uso y acceso a funcionalidades específicas de dispositivos particulares.

\subsection{\textit{BrainFlow}: plataforma unificada de adquisición}
\label{subsec:brainflow}

\textit{BrainFlow}~\cite{brainflow} ha emergido como la librería de abstracción más completa y ampliamente adoptada en el ecosistema \acrshort{bci} actual. Su desarrollo ha estado guiado por el objetivo de proporcionar una \acrshort{api} consistente que abstraiga las complejidades específicas de más de 20 fabricantes diferentes de dispositivos \acrshort{eeg}, desde sistemas clínicos de alta gama hasta dispositivos de consumo.

Sus características principales incluyen una abstracción de hardware que permite el acceso uniforme a dispositivos de diferentes fabricantes mediante una \acrshort{api} común, eliminando la necesidad de aprender interfaces específicas para cada tipo de hardware. El soporte multiplataforma abarca \textit{Windows}, \textit{macOS}, \textit{Linux} y \textit{Android}, facilitando el desarrollo de aplicaciones que puedan ejecutarse en diferentes entornos. La disponibilidad de \textit{bindings} para múltiples lenguajes de programación, incluyendo \textit{Python}, \textit{C++}, \textit{C\#}, \textit{Java}, \textit{R} y \textit{MATLAB}, amplía significativamente su accesibilidad para desarrolladores con diferentes antecedentes técnicos.

La librería proporciona capacidades de \textit{streaming} en tiempo real con latencias típicamente inferiores a 50 milisegundos, crucial para aplicaciones \acrshort{bci} interactivas. Adicionalmente, incluye implementaciones optimizadas de algoritmos básicos de filtrado digital que pueden aplicarse directamente a los datos adquiridos, reduciendo la complejidad del procesamiento posterior.

Sin embargo, \textit{BrainFlow} presenta limitaciones significativas que limitan su aplicabilidad como solución completa. La ausencia de interfaces gráficas de usuario requiere que los usuarios desarrollen sus propias herramientas de visualización y control. Las funcionalidades de procesamiento se limitan a operaciones básicas, requiriendo librerías adicionales para análisis avanzados. La falta de herramientas integradas para sincronización con eventos externos complica el desarrollo de aplicaciones que requieren correlación precisa entre señales cerebrales y acciones del usuario.

\subsection{\textit{Lab Streaming Layer}: protocolo de sincronización}
\label{subsec:lsl}

\textit{Lab Streaming Layer} (\acrshort{lsl})~\cite{kothe2013lab} representa un enfoque diferente al problema de integración en sistemas \acrshort{bci}, centrándose específicamente en el \textit{streaming} y sincronización de datos temporales en experimentos de investigación. Su diseño está optimizado para escenarios donde múltiples fuentes de datos deben combinarse con sincronización temporal precisa.

Sus características principales incluyen un protocolo de red que permite el descubrimiento automático de \textit{streams} de datos en la red local, facilitando la configuración de sistemas distribuidos. El sistema de sincronización temporal incorpora corrección automática de deriva de reloj entre diferentes dispositivos, crucial para mantener la precisión temporal en grabaciones largas. La capacidad de manejar múltiples tipos de datos simultáneamente, incluyendo señales \acrshort{eeg}, marcadores de eventos, audio y video, lo convierte en una solución integral para experimentos multimodales. La robustez del sistema incluye tolerancia a interrupciones de red y capacidades de reconexión automática.

Por otra parte, las limitaciones de \acrshort{lsl} incluyen la necesidad de programación personalizada sustancial para implementar aplicaciones completas, ya que \acrshort{lsl} proporciona únicamente la infraestructura de \textit{streaming} sin interfaces de usuario o herramientas de análisis. La dependencia frecuente de software propietario del fabricante para acceder a muchos dispositivos comerciales complica la configuración y limita la portabilidad. La ausencia de herramientas de visualización integradas requiere el desarrollo de soluciones personalizadas para monitoreo en tiempo real. La curva de aprendizaje pronunciada puede ser intimidante para usuarios sin experiencia técnica avanzada.

\section{Software de procesamiento de señales \acrshort{eeg}}
\label{sec:software_procesamiento}

El procesamiento de señales \acrshort{eeg} requiere herramientas especializadas capaces de manejar las características únicas de estos datos biomédicos, incluyendo la presencia inevitable de artefactos, la necesidad de análisis en múltiples dominios temporales y frecuenciales, y la importancia crítica de la preservación de información temporal precisa. El ecosistema de software disponible refleja diferentes filosofías de diseño y objetivos de aplicación, desde herramientas de análisis offline hasta plataformas de desarrollo de aplicaciones en tiempo real.

La diversidad de herramientas disponibles presenta tanto oportunidades como desafíos para investigadores y desarrolladores. Mientras que la variedad permite seleccionar herramientas optimizadas para tareas específicas, también requiere inversiones significativas en aprendizaje y puede resultar en flujos de trabajo fragmentados cuando se requiere combinar múltiples herramientas para completar análisis complejos.

\subsection{Entornos de análisis científico}
\label{subsec:entornos_analisis}

Los entornos de análisis científico representan plataformas completas diseñadas para cubrir todo el flujo de trabajo de análisis de datos \acrshort{eeg}, desde el preprocesamiento básico hasta análisis estadísticos sofisticados y visualización de resultados. Estas herramientas han sido desarrolladas por y para la comunidad de neurociencia, incorporando décadas de conocimiento especializado en procesamiento de señales biomédicas.

\subsubsection{\textit{EEGLAB}: plataforma \textit{MATLAB} establecida}

\textit{EEGLAB}~\cite{delorme2004eeglab} ha establecido durante más de dos décadas el estándar de facto para análisis de datos \acrshort{eeg} en muchos laboratorios de neurociencia mundial. Su desarrollo ha estado guiado por la necesidad de proporcionar herramientas accesibles pero potentes para investigadores sin formación especializada en procesamiento de señales.

La plataforma proporciona un flujo de trabajo completo que abarca desde operaciones básicas de preprocesamiento, incluyendo filtrado, eliminación de artefactos y re-referenciación, hasta análisis estadísticos avanzados y comparaciones entre condiciones experimentales. Su implementación del \acrshort{ica} se considera una de las más robustas disponibles, permitiendo la separación efectiva de fuentes de artefactos y señales de interés. Las capacidades de visualización incluyen topografías de actividad cerebral, espectrogramas tiempo-frecuencia y representaciones tridimensionales de fuentes estimadas.

El ecosistema de \textit{plugins} de \textit{EEGLAB} es particularmente notable, con cientos de extensiones desarrolladas por la comunidad que abordan necesidades específicas de análisis. Esta extensibilidad ha permitido que \textit{EEGLAB} se adapte a nuevas metodologías y paradigmas experimentales sin requerir modificaciones del núcleo de la plataforma.

Sin embargo, las limitaciones de \textit{EEGLAB} incluyen su dependencia completa de \textit{MATLAB}, cuya licencia comercial costosa puede ser prohibitiva para muchas instituciones. Su orientación hacia análisis \textit{offline} lo hace inadecuado para aplicaciones de tiempo real. La interfaz de usuario, aunque funcional, puede resultar intimidante para usuarios principiantes, y el manejo de \textit{datasets} grandes puede verse limitado por las restricciones de memoria de \textit{MATLAB}.

\subsubsection{\textit{MNE-Python}: plataforma moderna de código abierto}

\textit{MNE-Python}~\cite{gramfort2013meg} representa la evolución natural de las herramientas de análisis \acrshort{eeg} hacia el ecosistema \textit{Python}, combinando las lecciones aprendidas de plataformas establecidas con el diseño moderno orientado a objetos y las ventajas del desarrollo de código abierto.

Su \acrshort{api} moderna se caracteriza por un diseño orientado a objetos intuitivo que refleja conceptos naturales del dominio \acrshort{eeg}, como experimentos, épocas y eventos. La documentación exhaustiva incluye tutoriales, ejemplos prácticos y referencias técnicas detalladas que facilitan tanto el aprendizaje inicial como el uso avanzado. El soporte nativo para formatos estándar de la industria, incluyendo \acrshort{bids}, \acrshort{edf}+ y \textit{FIF}, facilita la interoperabilidad con otras herramientas y el intercambio de datos entre laboratorios.

La integración profunda con el ecosistema científico de \textit{Python} permite aprovechar bibliotecas especializadas como \textit{scikit-learn} para \textit{machine learning}, \textit{scipy} para procesamiento de señales avanzado y \textit{matplotlib} para visualización personalizada. Las capacidades de visualización incluyen gráficos interactivos que permiten exploración dinámica de datos complejos.

El ecosistema \textit{MNE} se ha expandido para incluir \textit{MNE-MATLAB}~\cite{mne_matlab} para usuarios que prefieren permanecer en el entorno \textit{MATLAB}, \textit{MNE-C}~\cite{mne_c} para aplicaciones que requieren máximo rendimiento, y diversas herramientas especializadas para casos de uso específicos.

Las fortalezas particulares de \textit{MNE-Python} incluyen su integración natural con el ecosistema \textit{Python} que facilita el desarrollo de flujos de trabajo complejos, una comunidad activa que garantiza desarrollo continuo y soporte responsivo, flexibilidad que permite tanto análisis básicos como investigación metodológica avanzada, y operaciones vectorizadas optimizadas que proporcionan rendimiento competitivo con implementaciones compiladas.

\begin{table}[hp!]
\centering
\resizebox{\textwidth}{!}{
\begin{tabular}{l|c|c|c|c|c}
\rowcolor{udcpink!25}
\textbf{Plataforma} & \textbf{Lenguaje} & \textbf{Licencia} & \textbf{Tiempo Real} & \textbf{Formatos} & \textbf{Comunidad} \\
\hline
\textbf{EEGLAB} & MATLAB & GPL & No & Múltiples & Madura \\
\textbf{MNE-Python} & Python & BSD & Limitado & BIDS/EDF+/FIF & Activa \\
\textbf{FieldTrip} & MATLAB & GPL & No & Múltiples & Especializada \\
\textbf{BrainStorm} & MATLAB & GPL & No & Múltiples & Médica \\
\end{tabular}
}
\caption{Comparativa de entornos de análisis científico \acrshort{eeg}}
\label{tab:entornos_analisis}
\end{table}

\subsection{Herramientas de desarrollo específicas}
\label{subsec:herramientas_desarrollo}

Las herramientas de desarrollo específicas se distinguen de los entornos de análisis por su enfoque en la creación de aplicaciones \acrshort{bci} interactivas y sistemas de tiempo real, más que en el análisis posterior de datos experimentales. Estas plataformas incorporan capacidades de prototipado rápido, desarrollo visual y despliegue de aplicaciones que pueden funcionar en entornos de producción.

\textit{OpenViBE}~\cite{openvibe} se ha establecido como una plataforma líder para el desarrollo visual de aplicaciones \acrshort{bci}, proporcionando un entorno de desarrollo basado en cajas de procesamiento interconectadas que permite el diseño intuitivo de flujos de procesamiento complejos. Su capacidad de operar tanto en tiempo real como en modo \textit{offline} lo hace versátil para diferentes fases del desarrollo de aplicaciones. La interfaz modular facilita la experimentación con diferentes algoritmos y configuraciones sin requerir programación explícita.

Sin embargo, \textit{OpenViBE} está limitado principalmente a actividades de investigación y prototipado, con capacidades limitadas para el desarrollo de productos comerciales. La curva de aprendizaje puede ser significativa para usuarios sin experiencia previa en procesamiento de señales.

\textit{BCI2000}~\cite{schalk2004bci2000} representa un \textit{framework} más técnico desarrollado en \textit{C++} que proporciona una arquitectura modular robusta para adquisición, procesamiento y aplicación de señales \acrshort{bci}. Su soporte para múltiples paradigmas \acrshort{bci} lo hace versátil para diferentes tipos de aplicaciones. Sin embargo, su interfaz de configuración compleja y la pronunciada curva de aprendizaje limitan su accesibilidad para usuarios no especializados.

\subsection{Librerías especializadas complementarias}
\label{subsec:librerias_especializadas}

El ecosistema de herramientas \acrshort{eeg} incluye numerosas librerías especializadas que abordan necesidades específicas no cubiertas completamente por las plataformas principales. Estas herramientas representan innovaciones dirigidas por la comunidad que exploran nuevos enfoques metodológicos o atienden nichos particulares de aplicación.

\textit{NeuroKit2}~\cite{neurokit2} se ha posicionado como un \textit{framework} \textit{Python} integral para análisis de señales fisiológicas, extracción de características y visualización de datos neurofisiológicos. Su diseño modular y documentación accesible lo han convertido en una opción popular para investigadores que requieren análisis que van más allá del \acrshort{eeg} tradicional.

\textit{pyEEG}~\cite{pyeeg} se especializa en procesamiento continuo de datos \acrshort{eeg} y se basa en \textit{MNE-Python} y \textit{scikit-learn} para proporcionar capacidades avanzadas de generación de funciones de respuesta temporal. Su enfoque en señales continuas lo hace particularmente útil para aplicaciones de monitoreo a largo plazo.

\textit{EEG-ExPy}~\cite{eeg_expy} representa más una colección de experimentos en \textit{Python} que una librería tradicional, pero su compatibilidad con múltiples dispositivos comerciales (incluyendo \textit{Muse}, \textit{OpenBCI}, \textit{Neurosity} y \textit{g.tec}) lo convierte en un recurso valioso para prototipado rápido de experimentos.

\begin{table}[hp!]
\centering
\resizebox{\textwidth}{!}{
\begin{tabular}{l|c|c|c|c}
\rowcolor{udcpink!25}
\textbf{Librería} & \textbf{Enfoque} & \textbf{Dispositivos} & \textbf{Mantenimiento} & \textbf{Documentación} \\
\hline
\textbf{NeuroKit2} & Señales fisiológicas & Genérico & Activo & Excelente \\
\textbf{pyEEG} & Procesamiento continuo & MNE-compatible & Moderado & Buena \\
\textbf{EEG-ExPy} & Experimentos & Múltiples & Activo & Básica \\
\textbf{Emokit} & Emotiv específico & Emotiv EPOC & Discontinuado & Limitada \\
\end{tabular}
}
\caption{Comparativa de librerías especializadas complementarias}
\label{tab:librerias_especializadas}
\end{table}

\section{Análisis de limitaciones y oportunidades}
\label{sec:analisis_limitaciones}

El análisis exhaustivo del ecosistema actual de herramientas \acrshort{bci} revela un patrón consistente de fragmentación que, aunque refleja la diversidad y vitalidad del campo, también crea barreras significativas para la adopción práctica de tecnologías \acrshort{bci}. Esta fragmentación se manifiesta en múltiples dimensiones que van desde la incompatibilidad técnica hasta las barreras de conocimiento especializado.

La fragmentación del ecosistema constituye quizás la limitación más fundamental identificada. Las herramientas existentes abordan aspectos específicos del desarrollo \acrshort{bci} de forma aislada, requiriendo que los desarrolladores combinen múltiples soluciones para crear aplicaciones completas. \textit{BrainFlow} proporciona capacidades excelentes para adquisición de datos pero carece de herramientas integradas de procesamiento avanzado y visualización. \textit{MNE-Python} ofrece análisis sofisticado pero no está diseñado para operación en tiempo real. \textit{LSL} facilita la sincronización de datos pero requiere programación extensiva para implementar funcionalidades de aplicación.

La ausencia de herramientas para captura sincronizada representa una limitación crítica específica. Ninguna de las soluciones analizadas proporciona capacidades integradas para capturar simultáneamente señales \acrshort{eeg} y eventos de entrada del usuario con la sincronización temporal precisa requerida para aplicaciones \acrshort{bci} efectivas. Esta carencia obliga a los desarrolladores a implementar soluciones personalizadas que son propensas a errores y difíciles de validar.

Las barreras de entrada técnicas constituyen otro obstáculo significativo. La mayoría de herramientas disponibles asumen conocimientos avanzados de procesamiento de señales, neurociencia y programación, creando barreras prohibitivas para investigadores en disciplinas relacionadas que podrían beneficiarse de tecnologías \acrshort{bci}. Esta situación limita la diversidad de perspectivas y aplicaciones que podrían enriquecer el campo.

La carencia de interfaces de usuario intuitivas agrava las barreras de entrada. Pocas herramientas proporcionan interfaces gráficas que permitan a usuarios no técnicos explorar y utilizar capacidades \acrshort{bci} de forma productiva. Esta limitación restringe el campo principalmente a especialistas técnicos, limitando la exploración de aplicaciones novedosas que podrían emerger de disciplinas diversas.

La falta de estandarización en formatos de datos complica la interoperabilidad entre herramientas y la reproducibilidad de investigación. Cada herramienta tiende a utilizar formatos propietarios o específicos que requieren conversiones complejas y propensas a errores cuando se desea combinar datos de diferentes fuentes o migrar entre herramientas.

Sin embargo, estas limitaciones revelan oportunidades significativas para el desarrollo de soluciones más integradas y accesibles. La necesidad evidente de una plataforma integrada que combine adquisición, procesamiento y sincronización en una solución cohesiva representa una oportunidad clara para agregar valor significativo al ecosistema. La demanda de interfaces múltiples que soporten tanto usuarios técnicos como no técnicos puede democratizar el acceso a tecnologías \acrshort{bci} y facilitar la exploración de aplicaciones novedosas.

La carencia de soluciones que combinen análisis avanzado con capacidades de tiempo real representa otra oportunidad importante. El desarrollo de herramientas que mantengan la sofisticación de plataformas como \textit{MNE-Python} mientras proporcionan las capacidades de tiempo real necesarias para aplicaciones interactivas podría transformar las posibilidades de desarrollo \acrshort{bci}.

La necesidad identificada de herramientas que faciliten la generación de \textit{datasets} bien estructurados para \acrshort{bci} representa una oportunidad estratégica. El desarrollo de soluciones que automaticen  prácticas para captura, organización y documentación de datos podría acelerar significativamente el progreso en investigación \acrshort{bci} y el desarrollo de aplicaciones prácticas.

La relación entre estas limitaciones y oportunidades sugiere que el desarrollo de soluciones integradas no solo es técnicamente factible sino estratégicamente valioso. La convergencia de tecnologías maduras en procesamiento de señales, desarrollo de interfaces y distribución de software crea un contexto favorable para el desarrollo de plataformas que puedan abordar múltiples limitaciones simultáneamente mientras proporcionan valor agregado significativo a la comunidad \acrshort{bci}.

\section{Posicionamiento estratégico de la solución propuesta}
\label{sec:posicionamiento_estrategico}

El análisis del panorama actual de herramientas \acrshort{bci} revela un espacio de oportunidad claro para una solución que integre las fortalezas de las herramientas existentes mientras aborda sus limitaciones fundamentales. La solución propuesta en este trabajo se posiciona estratégicamente para llenar este vacío mediante una aproximación que combina integración técnica, accesibilidad de interfaces y extensibilidad arquitectónica.

La diferenciación principal de la solución propuesta reside en su enfoque integral hacia el problema de fragmentación del ecosistema. Mientras que las herramientas existentes optimizan aspectos específicos del desarrollo \acrshort{bci}, la plataforma desarrollada proporciona capacidades integradas que abarcan todo el flujo de trabajo desde la adquisición hasta el análisis, eliminando la necesidad de combinar múltiples herramientas manualmente.

La estrategia de múltiples interfaces permite que la misma plataforma sirva a usuarios con diferentes niveles técnicos y contextos de uso. Esta aproximación contrasta con la tendencia predominante en el ecosistema actual hacia herramientas especializadas que atienden segmentos específicos de usuarios, proporcionando en su lugar una solución unificada que puede evolucionar con las necesidades cambiantes de los usuarios.

La arquitectura extensible basada en principios \acrshort{solid} y patrones dirigidos por eventos posiciona la solución para evolucionar con los avances tecnológicos del campo \acrshort{bci}, contrastando con enfoques más rígidos que requieren modificaciones fundamentales para incorporar nuevas capacidades o dispositivos.

El enfoque en la generación de \textit{datasets} de calidad responde a una necesidad fundamental del campo que no está siendo atendida adecuadamente por las herramientas existentes. Esta capacidad posiciona la solución como una herramienta valiosa no solo para el desarrollo de aplicaciones sino también para la investigación fundamental en \acrshort{bci}.

\section{Resumen del análisis del estado del arte}
\label{sec:resumen_estado_arte}

El ecosistema actual de herramientas para desarrollo de sistemas \acrshort{bci} presenta una diversidad rica que refleja tanto la maduración del campo como su fragmentación inherente. Mientras que existen soluciones robustas para aspectos específicos del desarrollo \acrshort{bci}, incluyendo adquisición de datos (\textit{BrainFlow}), análisis avanzado (\textit{MNE-Python}) y sincronización de datos (\textit{LSL}), ninguna proporciona una solución integrada que combine estas capacidades con interfaces accesibles y herramientas específicas para generación de \textit{datasets} de calidad.

Esta fragmentación, aunque comprensible dada la evolución histórica del campo, crea barreras significativas para la adopción práctica de tecnologías \acrshort{bci} y limita las oportunidades de innovación interdisciplinaria. Las limitaciones identificadas incluyen la ausencia de herramientas para captura sincronizada, barreras de entrada técnicas elevadas, carencia de interfaces intuitivas y falta de estandarización en formatos de datos.

Sin embargo, estas limitaciones revelan oportunidades claras para el desarrollo de soluciones más integradas que puedan democratizar el acceso a tecnologías \acrshort{bci} mientras mantienen las capacidades avanzadas requeridas por usuarios especializados. El contexto tecnológico actual, caracterizado por la disponibilidad de herramientas de desarrollo maduras y estándares emergentes, proporciona una base favorable para el desarrollo de tales soluciones integradas.